\section{示例问答与测试用例}

下列为系统在真实 16 万篇语料上的\textbf{实际输出}(非示意),用于直观展示检索、跨语言、推荐与证据接地能力。

\subsection{用例一:英文混合检索(双臂命中)}
\begin{plainbox}
\textbf{查询}:\texttt{graph neural network materials discovery}\par\vspace{2pt}
\textbf{返回}(\statusbadge{lexical}+\statusbadge{semantic} 双臂命中):
\begin{enumerate}
\item Accelerated Discovery of Energy Materials via Graph Neural Networks \hfill(2025)
\item CTGNN: Crystal Transformer Graph Neural Network for Crystal Materials \hfill(2024)
\item DeepCrysTet: A Deep Learning Approach Using Tetrahedral Mesh \hfill(2023)
\end{enumerate}
\end{plainbox}

\subsection{用例二:中文跨语言检索(命中英文文献)}
\begin{plainbox}
\textbf{查询}:\texttt{大语言模型的检索增强生成}(中文)\par\vspace{2pt}
\textbf{返回}(\statusbadge{semantic} 语义臂,跨语言命中英文论文):
\begin{enumerate}
\item Retrieval-Augmented Retrieval: Large Language Models are Strong Zero-Shot Retrievers \hfill(2024)
\item LExT: Towards Evaluating Trustworthiness of Natural Language Explanations \hfill(2025)
\item A Comprehensive Study of Jailbreak Attack versus Defense for LLMs \hfill(2024)
\end{enumerate}
\end{plainbox}

\subsection{用例三:可解释论文推荐}
\begin{plainbox}
\textbf{种子论文}:\texttt{PMC13273979}(AI 药物发现方向)\par\vspace{2pt}
\textbf{推荐}(每条附可解释因子 semantic / keyword / author / recency):
\begin{center}
\begin{tabularx}{\textwidth}{Xccc}
\toprule
\textbf{推荐论文} & \textbf{年份} & \textbf{语义相似} & \textbf{融合因子} \\
\midrule
Generative AI in drug discovery and development & 2024 & 0.941 & sem .56 / kw .10 \\
Editorial: Advancing drug discovery with AI & 2025 & 0.934 & sem .56 / kw .10 \\
Advancing Drug Discovery with AI: Machine and Deep Learning & 2026 & 0.915 & sem .55 / kw .10 \\
\bottomrule
\end{tabularx}
\end{center}
\end{plainbox}

\subsection{用例四:趋势热点(模型输出)}
\begin{plainbox}
\textbf{\texttt{/api/trends} 新兴热点排行}(关键词清洗后):
\texttt{retrieval augmented generation} → \texttt{large language model} → \texttt{gene delivery} → \dots\par\vspace{2pt}
与 2022--2026 年科研态势吻合,RAG 与大语言模型领跑。
\end{plainbox}

\subsection{用例五:证据接地问答}
\begin{plainbox}
\textbf{提问}:图神经网络在材料发现中有哪些应用?\par\vspace{2pt}
\textbf{系统流程}:混合检索 → 取 Top 证据(如 \emph{A review on the applications of GNNs in materials science}、\emph{crystal graph convolutional neural network})→ 本地大模型\textbf{仅基于证据}生成回答,并标注引用 [n]。\par\vspace{2pt}
\textbf{关键特性}:回答可追溯到具体论文;证据不足时如实说明,不臆造。
\end{plainbox}
